\documentclass[12pt]{article}

% Ћирилица: кодирање фонта и језик (српски, ћирилица)
\usepackage[T2A]{fontenc}
\usepackage[utf8]{inputenc}
\usepackage[serbianc]{babel}

% Величина стране и маргине
\usepackage[a4paper,top=2cm,bottom=2cm,left=2.5cm,right=2.5cm]{geometry}

% Корисни пакети
\usepackage{amsmath}
\usepackage{enumitem}
\usepackage{tcolorbox}
\tcbuselibrary{breakable}
\usepackage[colorlinks=true, allcolors=blue]{hyperref}

% Цртање електричних кола (европски симболи)
\usepackage[european]{circuitikz}
\usetikzlibrary{arrows.meta, decorations.markings, babel}

% Стрелица на средини линије (за струје у чвору)
\tikzset{
    sredstrelica/.style={postaction={decorate}, decoration={markings,
        mark=at position 0.55 with {\arrow{Stealth[length=7pt]}}}},
    obilazak/.style={-{Stealth[length=6pt]}, gray!70!black},
}

% Оквир за задатке
% (без breakable: задатак се не дели на две стране)
\newtcolorbox[auto counter]{zadatak}{
    colback=white, colframe=black!55, colbacktitle=black!55, coltitle=white,
    fonttitle=\bfseries\small, title={Задатак~\thetcbcounter},
    boxrule=0.5pt, arc=1.5pt, left=6pt, right=6pt, top=3pt, bottom=3pt,
    before skip=10pt, after skip=12pt}

% Оквир за решења
\newtcolorbox{resenje}[1]{breakable,
    colback=blue!3, colframe=blue!45!black, colbacktitle=blue!45!black, coltitle=white,
    fonttitle=\bfseries\small, title={Решење задатка~#1},
    boxrule=0.5pt, arc=1.5pt, left=6pt, right=6pt, top=3pt, bottom=3pt,
    before skip=10pt, after skip=12pt}

\title{Кирхофови закони -- задаци}
\author{Основи електротехнике 1}
\date{}

\begin{document}
\maketitle

\noindent Решења задатака налазе се на крају документа.

%==================================================================
% ЗАДАЦИ
%==================================================================

% ---------- Задатак 1: лак, први Кирхофов закон ----------
\begin{zadatak}
\textit{(Први Кирхофов закон)}

У чвор на слици улази струја $I = 10\,\text{A}$. Из чвора излазе струје
$I_1 = 3\,\text{A}$, $I_2 = 5\,\text{A}$ и струја $I_3$. Израчунати струју $I_3$.

\begin{center}
\begin{tikzpicture}[thick]
    \draw[sredstrelica] (-2.2,0) -- (0,0) node[pos=0.45, above]{$I$};
    \draw[sredstrelica] (0,0)    -- (0,1.8) node[pos=0.55, right]{$I_1$};
    \draw[sredstrelica] (0,0)    -- (2.2,0) node[pos=0.55, above]{$I_2$};
    \draw[sredstrelica] (0,0)    -- (0,-1.8) node[pos=0.55, right]{$I_3$};
    \fill (0,0) circle (3pt);
\end{tikzpicture}
\end{center}
\end{zadatak}

% ---------- Задатак 2: лак, други Кирхофов закон ----------
\begin{zadatak}
\textit{(Други Кирхофов закон)}

У колу на слици је $E = 18\,\text{V}$, $R_1 = 2\,\Omega$, $R_2 = 4\,\Omega$ и
$R_3 = 3\,\Omega$. Израчунати струју у колу и напоне на сва три отпорника.

\begin{center}
\begin{circuitikz}
    \draw (0,0) to[battery1, l=$E$, invert] (0,3)
          to[short, i>^=$I$] (1,3)
          to[R, l=$R_1$] (4,3)
          to[R, l=$R_2$] (4,0)
          to[R, l=$R_3$] (0,0);
\end{circuitikz}
\end{center}
\end{zadatak}

% ---------- Задатак 3: тежи, први Кирхофов закон ----------
\begin{zadatak}
\textit{(Први Кирхофов закон)}

На слици је приказан део кола са два чвора, $A$ и $B$. Познато је
$I_1 = 8\,\text{A}$, $I_2 = 3\,\text{A}$ и $I_4 = 9\,\text{A}$. Смерови струја
означени на слици су \textbf{претпостављени}.
\begin{enumerate}
    \item[а)] Израчунати струју $I_3$.
    \item[б)] Израчунати струју $I_5$ и одредити њен стварни смер.
\end{enumerate}

\begin{center}
\begin{tikzpicture}[thick]
    \draw[sredstrelica] (-2.2,0) -- (0,0) node[pos=0.45, above]{$I_1$};
    \draw[sredstrelica] (0,0)    -- (0,-1.8) node[pos=0.55, left]{$I_2$};
    \draw[sredstrelica] (0,0)    -- (3.5,0) node[pos=0.5, above]{$I_3$};
    \draw[sredstrelica] (3.5,0)  -- (5.7,0) node[pos=0.55, above]{$I_4$};
    \draw[sredstrelica] (3.5,0)  -- (3.5,-1.8) node[pos=0.55, right]{$I_5$};
    \fill (0,0) circle (3pt) node[above left]{$A$};
    \fill (3.5,0) circle (3pt) node[above right]{$B$};
\end{tikzpicture}
\end{center}
\end{zadatak}

% ---------- Задатак 4: тежи, други Кирхофов закон ----------
\begin{zadatak}
\textit{(Други Кирхофов закон)}

У колу на слици је $E_1 = 30\,\text{V}$, $E_2 = 6\,\text{V}$,
$R_1 = 3\,\Omega$ и $R_2 = 5\,\Omega$. Струја у колу је $I = 2\,\text{A}$ и
тече у смеру означеном на слици.
\begin{enumerate}
    \item[а)] Израчунати отпорност $R_3$.
    \item[б)] Израчунати напоне на сва три отпорника.
    \item[в)] Проверити резултат помоћу другог Кирхофовог закона.
\end{enumerate}

\begin{center}
\begin{circuitikz}
    \draw (0,0) to[battery1, l=$E_1$, invert] (0,3)
          to[short, i>^=$I$] (0.8,3)
          to[R, l=$R_1$] (3,3)
          to[R, l=$R_2$] (5.5,3)
          to[battery1, l=$E_2$] (5.5,0)
          to[R, l=$R_3$] (0,0);
\end{circuitikz}
\end{center}
\end{zadatak}

% ---------- Задатак 5: комбинација оба закона ----------
\begin{zadatak}
\textit{(Први и други Кирхофов закон)}

У колу на слици је $E = 24\,\text{V}$, $R_1 = 4\,\Omega$, $R_2 = 6\,\Omega$ и
$R_3 = 12\,\Omega$. Израчунати:
\begin{enumerate}
    \item[а)] еквивалентну отпорност кола,
    \item[б)] струју $I$ коју даје извор,
    \item[в)] напон на отпорнику $R_1$ и напон између чворова $A$ и $B$,
    \item[г)] струје $I_2$ и $I_3$,
    \item[д)] проверити резултате помоћу првог и другог Кирхофовог закона.
\end{enumerate}

\begin{center}
\begin{circuitikz}
    \draw (0,0) to[battery1, l=$E$, invert] (0,3)
          to[short, i>^=$I$] (0.8,3)
          to[R, l=$R_1$] (3.2,3)
          to[short, -*] (4,3)
          to[R, l_=$R_2$, i>^=$I_2$, -*] (4,0)
          -- (0,0);
    \draw (4,3) -- (7,3) to[R, l_=$R_3$, i>^=$I_3$] (7,0) -- (4,0);
    \node[above] at (4,3) {$A$};
    \node[below] at (4,0) {$B$};
\end{circuitikz}
\end{center}
\end{zadatak}

%==================================================================
% РЕШЕЊА
%==================================================================
\newpage
\section*{Решења}

% ---------- Решење 1 ----------
\begin{resenje}{1}
Струја $I$ улази у чвор, а струје $I_1$, $I_2$ и $I_3$ излазе из њега:
\[
I = I_1 + I_2 + I_3
\quad\Longrightarrow\quad
I_3 = I - I_1 - I_2 = 10\,\text{A} - 3\,\text{A} - 5\,\text{A} = 2\,\text{A}
\]
\end{resenje}

% ---------- Решење 2 ----------
\begin{resenje}{2}
Обилазимо контуру у смеру казаљке на сату, у смеру струје. Кроз извор
пролазимо од $-$ ка $+$, па је $E$ позитивна:
\[
E = R_1 I + R_2 I + R_3 I = (R_1 + R_2 + R_3)\, I
\]
\[
I = \frac{E}{R_1 + R_2 + R_3}
  = \frac{18\,\text{V}}{2\,\Omega + 4\,\Omega + 3\,\Omega}
  = \frac{18\,\text{V}}{9\,\Omega} = 2\,\text{A}
\]
Напони на отпорницима:
\[
U_1 = R_1 I = 4\,\text{V}, \qquad
U_2 = R_2 I = 8\,\text{V}, \qquad
U_3 = R_3 I = 6\,\text{V}
\]
\textbf{Провера:} $U_1 + U_2 + U_3 = 4\,\text{V} + 8\,\text{V} + 6\,\text{V}
= 18\,\text{V} = E$.
\end{resenje}

% ---------- Решење 3 ----------
\begin{resenje}{3}
а) \textbf{Чвор $A$:} струја $I_1$ улази, а $I_2$ и $I_3$ излазе:
\[
I_1 = I_2 + I_3
\quad\Longrightarrow\quad
I_3 = I_1 - I_2 = 8\,\text{A} - 3\,\text{A} = 5\,\text{A}
\]

б) \textbf{Чвор $B$:} струја $I_3$ улази, а $I_4$ и $I_5$ (по претпоставци)
излазе:
\[
I_3 = I_4 + I_5
\quad\Longrightarrow\quad
I_5 = I_3 - I_4 = 5\,\text{A} - 9\,\text{A} = -4\,\text{A}
\]
Резултат је негативан, па струја $I_5$ у ствари \textbf{улази} у чвор $B$ и
има јачину $4\,\text{A}$.

\textbf{Провера за чвор $B$:} улази $5\,\text{A} + 4\,\text{A} = 9\,\text{A}$,
излази $9\,\text{A}$.
\end{resenje}

% ---------- Решење 4 ----------
\begin{resenje}{4}
а) Обилазимо контуру у смеру казаљке на сату, у смеру струје:
\begin{itemize}
    \item кроз $E_1$ пролазимо од $-$ ка $+$, па је $E_1$ позитивна,
    \item кроз $E_2$ пролазимо од $+$ ка $-$, па је $E_2$ негативна,
    \item струја тече у смеру обиласка, па су сви падови напона позитивни.
\end{itemize}
\[
E_1 - E_2 = (R_1 + R_2 + R_3)\, I
\]
Одавде изражавамо $R_3$:
\[
R_1 + R_2 + R_3 = \frac{E_1 - E_2}{I} = \frac{30\,\text{V} - 6\,\text{V}}{2\,\text{A}}
= 12\,\Omega
\]
\[
R_3 = 12\,\Omega - R_1 - R_2 = 12\,\Omega - 3\,\Omega - 5\,\Omega = 4\,\Omega
\]

б) Напони на отпорницима:
\[
U_1 = R_1 I = 6\,\text{V}, \qquad
U_2 = R_2 I = 10\,\text{V}, \qquad
U_3 = R_3 I = 8\,\text{V}
\]

в) \textbf{Провера:}
\[
U_1 + U_2 + U_3 = 6\,\text{V} + 10\,\text{V} + 8\,\text{V} = 24\,\text{V},
\qquad
E_1 - E_2 = 30\,\text{V} - 6\,\text{V} = 24\,\text{V}
\]
\end{resenje}

% ---------- Решење 5 ----------
\begin{resenje}{5}
а) Отпорници $R_2$ и $R_3$ су паралелно:
\[
R_{23} = \frac{R_2 \cdot R_3}{R_2 + R_3}
       = \frac{6\,\Omega \cdot 12\,\Omega}{6\,\Omega + 12\,\Omega}
       = \frac{72}{18}\,\Omega = 4\,\Omega
\]
Са њима је редно везан $R_1$:
\[
R_e = R_1 + R_{23} = 4\,\Omega + 4\,\Omega = 8\,\Omega
\]

б) Струја извора:
\[
I = \frac{E}{R_e} = \frac{24\,\text{V}}{8\,\Omega} = 3\,\text{A}
\]

в) Напон на $R_1$ и напон између чворова $A$ и $B$:
\[
U_1 = R_1 I = 4\,\Omega \cdot 3\,\text{A} = 12\,\text{V}, \qquad
U_{AB} = R_{23} I = 4\,\Omega \cdot 3\,\text{A} = 12\,\text{V}
\]

г) Отпорници $R_2$ и $R_3$ су оба прикључени између чворова $A$ и $B$, па је
на оба напон $U_{AB}$:
\[
I_2 = \frac{U_{AB}}{R_2} = \frac{12\,\text{V}}{6\,\Omega} = 2\,\text{A}, \qquad
I_3 = \frac{U_{AB}}{R_3} = \frac{12\,\text{V}}{12\,\Omega} = 1\,\text{A}
\]

д) \textbf{Провера првим законом (чвор $A$):} струја $I$ улази, а $I_2$ и
$I_3$ излазе:
\[
I_2 + I_3 = 2\,\text{A} + 1\,\text{A} = 3\,\text{A} = I
\]

\textbf{Провера другим законом (контура $E \to R_1 \to R_2$):}
\[
R_1 I + R_2 I_2 = 4\,\Omega \cdot 3\,\text{A} + 6\,\Omega \cdot 2\,\text{A}
= 12\,\text{V} + 12\,\text{V} = 24\,\text{V} = E
\]
\end{resenje}

\end{document}